\chapter{Implémentation du prototype}
\label{chap:implementation}

\section{Introduction}

Le chapitre précédent a présenté Ariel \logminer{} comme une architecture modulaire
organisant plusieurs fonctions : collecte, parsing, normalisation, routage, détection,
corrélation, mémoire, audit, communication entre agents et restitution des résultats.
Le présent chapitre montre comment ces principes sont traduits dans l'implémentation
logicielle du prototype.

L'objectif n'est pas de reproduire le code source dans le mémoire. Une telle approche
rendrait le document difficile à lire et apporterait peu d'information scientifique.
L'implémentation est plutôt décrite à travers les responsabilités des modules, leurs
interfaces, les données qu'ils manipulent et les choix techniques qui permettent de relier
le modèle architectural à un système exécutable.

Le prototype est développé principalement en Python. Ce choix est lié à la disponibilité
d'un écosystème mature pour la manipulation de données, le machine learning, les API et
l'automatisation. Les modèles légers utilisés dans plusieurs expériences sont principalement
implémentés au moyen de scikit-learn. Redis intervient dans la distribution de certaines
tâches, tandis qu'une API locale et une interface web assurent l'accès aux fonctions de
consultation et de contrôle.

L'implémentation reste volontairement orientée vers un environnement local et de
laboratoire. Elle cherche d'abord à assurer la cohérence des traitements, leur
auditabilité et leur reproductibilité avant de considérer des propriétés telles que la
haute disponibilité ou le déploiement à grande échelle.

\section{Organisation générale du projet logiciel}

\subsection{Séparation entre code applicatif, données et expériences}

Le projet distingue plusieurs catégories de ressources.

Le code applicatif est regroupé autour du paquet Python
\texttt{src/logminer}. Il contient les composants nécessaires au pipeline,
aux agents, au routage, à l'audit et aux autres fonctions du prototype.

Les scripts d'expérimentation sont séparés du noyau applicatif. Ils servent notamment à
exécuter des campagnes contrôlées, comparer des modèles, mesurer les ressources, produire
des rapports et générer certains artefacts utilisés dans le mémoire.

Les données sont organisées selon leur rôle. Les sources brutes sont séparées des résultats
issus des transformations et des campagnes expérimentales. Cette organisation limite le
risque d'écraser une donnée d'origine pendant un traitement et facilite la relation entre
un résultat scientifique et le fichier qui l'a produit.

Enfin, les figures, tableaux et rapports de synthèse sont considérés comme des artefacts
dérivés. Leur rôle est de présenter les résultats, non de remplacer les données de preuve.

Cette organisation peut être résumée par la chaîne suivante :

\begin{center}
\textit{données sources}
$\rightarrow$
\textit{code de traitement}
$\rightarrow$
\textit{artefacts expérimentaux}
$\rightarrow$
\textit{figures et tableaux}
$\rightarrow$
\textit{mémoire}.
\end{center}

L'intérêt de cette séparation est méthodologique. Une valeur citée dans le manuscrit doit,
autant que possible, pouvoir être reliée à un fichier de résultat, puis au script et aux
données ayant servi à la produire.

\subsection{Modules principaux}

Le tableau~\ref{tab:modules-logminer} résume les responsabilités des principaux modules
du prototype.

\begin{table}[H]
\centering
\caption{Organisation fonctionnelle des principaux modules d'Ariel \logminer{}}
\label{tab:modules-logminer}
\small
\begin{tabularx}{\textwidth}{p{4cm}X}
\toprule
\textbf{Module ou groupe de modules} & \textbf{Responsabilité} \\
\midrule

\texttt{pipeline.py} &
Coordination de la chaîne locale de traitement et enchaînement des principales étapes. \\

Parseurs et adaptateurs &
Lecture des différents formats et extraction des informations disponibles. \\

Normalisation &
Projection des champs disponibles vers une représentation commune lorsqu'un adaptateur
compatible existe. \\

\texttt{model\_router.py} &
Identification de la famille probable d'une source et sélection du traitement correspondant. \\

Registre et modèles &
Association entre familles de données, caractéristiques attendues et artefacts de modèles. \\

Corrélateur &
Regroupement d'anomalies candidates selon les règles et informations contextuelles
disponibles. \\

\texttt{agents/bus.py} &
Définition du contrat métier des messages et fonctions de communication associées. \\

\texttt{supervisor\_agent.py} et workers &
Sélection, distribution, suivi et exécution de certaines tâches par les agents. \\

Audit et mémoire &
Persistance des états, décisions, erreurs et informations nécessaires à la traçabilité. \\

API locale &
Exposition contrôlée des informations nécessaires au dashboard et à certains traitements. \\

Dashboard &
Consultation des anomalies candidates, incidents, ressources et informations d'audit. \\

\bottomrule
\end{tabularx}
\end{table}

Cette séparation n'implique pas que chaque module s'exécute obligatoirement sur une machine
distincte. Elle représente une décomposition logique. Dans certaines configurations,
plusieurs fonctions s'exécutent dans le même environnement local ; dans les campagnes
distribuées, certains traitements sont confiés à plusieurs workers.

\section{Environnement logiciel principal}

L'environnement expérimental final utilisé pour les campagnes documentées repose sur
Python 3.11.9. Les versions de plusieurs bibliothèques importantes ont été gelées ou
relevées afin de faciliter la reproductibilité.

Le tableau~\ref{tab:versions-implementation} présente les versions attestées dans
l'environnement expérimental final.

\begin{table}[H]
\centering
\caption{Principales versions logicielles documentées dans l'environnement final}
\label{tab:versions-implementation}
\small
\begin{tabular}{ll}
\toprule
\textbf{Composant} & \textbf{Version documentée} \\
\midrule
Python & 3.11.9 \\
scikit-learn & 1.7.2 \\
pandas & 2.3.2 \\
NumPy & 2.2.3 \\
SciPy & 1.16.1 \\
psutil & 7.2.2 \\
joblib & 1.5.2 \\
Client Python Redis & 8.0.1 \\
Drain3 & 0.9.11 \\
\bottomrule
\end{tabular}
\end{table}

Les versions qui ne sont pas établies par les artefacts finaux ne sont pas inventées.
Cette règle concerne notamment certains composants de service dont la version précise
n'est pas nécessaire à la compréhension de l'architecture.

Les modèles de machine learning légers sont sérialisés au moyen des mécanismes compatibles
avec l'écosystème scikit-learn. Les traitements tabulaires s'appuient principalement sur
pandas et NumPy. psutil est utilisé dans les campagnes nécessitant des mesures de
ressources telles que la consommation CPU ou la mémoire résidente.

\section{Implémentation du pipeline principal}

\subsection{Rôle du pipeline}

Le pipeline constitue la colonne vertébrale du traitement local. Sa fonction est
d'enchaîner les composants sans confondre leurs responsabilités.

Le cycle général est :

\begin{enumerate}
    \item recevoir ou découvrir une source ;
    \item déterminer les informations disponibles sur cette source ;
    \item appliquer le parseur correspondant lorsque celui-ci existe ;
    \item produire une représentation intermédiaire ;
    \item normaliser les champs compatibles ;
    \item déterminer une famille de traitement ;
    \item sélectionner un modèle ou un handler ;
    \item exécuter le traitement ;
    \item enregistrer les sorties ;
    \item transmettre les anomalies candidates aux fonctions aval lorsque nécessaire.
\end{enumerate}

Le pipeline ne suppose pas que toutes les étapes réussissent. Une source peut être
correctement lue mais ne pas disposer d'un adaptateur complet. De même, le routeur peut
produire un fallback lorsque les informations disponibles ne permettent pas une décision
suffisamment cohérente.

Cette approche évite de considérer l'absence d'exception logicielle comme une preuve de
traitement correct.

\subsection{Contrôle des entrées}

L'entrée du pipeline peut correspondre à un fichier individuel, un groupe de fichiers ou
une unité produite par un collecteur.

Avant d'appliquer un modèle, plusieurs contrôles sont nécessaires :

\begin{itemize}
    \item existence et lisibilité de la source ;
    \item présence des colonnes ou champs requis ;
    \item reconnaissance du format ;
    \item cohérence minimale des valeurs ;
    \item disponibilité du traitement attendu.
\end{itemize}

Lorsqu'un contrôle échoue, le comportement souhaité n'est pas de fabriquer une valeur
par défaut afin de poursuivre silencieusement. Le système doit enregistrer un état
permettant de distinguer une entrée correctement transformée d'une entrée partiellement
traitée ou rejetée.

\section{Parsing des journaux}

\subsection{Principe d'implémentation}

Les parseurs transforment les représentations propres aux sources en objets pouvant être
manipulés par la suite du pipeline.

Un parseur doit répondre à trois questions :

\begin{enumerate}
    \item reconnaît-il le format reçu ?
    \item quels champs peut-il extraire de manière fiable ?
    \item que doit-il faire lorsque l'entrée n'est que partiellement conforme ?
\end{enumerate}

Cette troisième question est importante. Dans une application de sécurité, supprimer une
ligne parce qu'elle ne correspond pas exactement au format prévu peut créer un angle mort.
À l'inverse, déclarer qu'une ligne est correctement parsée lorsque peu d'informations sont
réellement extraites fausse l'évaluation.

Le comportement retenu privilégie donc un statut explicite et un fallback lorsque
nécessaire.

\subsection{Sources structurées}

Les fichiers CSV et autres structures tabulaires permettent généralement d'utiliser
directement les colonnes disponibles.

Le traitement vérifie notamment :

\begin{itemize}
    \item les noms de colonnes ;
    \item la présence de la variable cible lorsqu'une expérience supervisée l'exige ;
    \item le type des valeurs ;
    \item les valeurs manquantes ;
    \item les colonnes non numériques lorsque le modèle attend exclusivement des
    caractéristiques numériques.
\end{itemize}

Pour les expériences CICIDS2017, le pipeline expérimental sélectionne les caractéristiques
numériques prévues par le protocole final. Les labels sont conservés séparément de
l'entrée du modèle afin d'éviter qu'ils participent accidentellement à la prédiction.

\subsection{Sources textuelles}

Les journaux Linux, Apache ou d'autres systèmes textuels nécessitent une extraction
progressive.

Selon la source, les informations recherchées peuvent inclure :

\begin{itemize}
    \item timestamp ;
    \item hôte ;
    \item service ;
    \item processus ;
    \item utilisateur ;
    \item adresse ;
    \item niveau ;
    \item message.
\end{itemize}

Tous ces champs ne sont pas disponibles dans toutes les sources. L'implémentation doit
donc accepter des événements partiellement structurés.

\subsection{Windows et exports structurés}

Les événements Windows possèdent une structure différente des journaux textuels simples.
L'implémentation peut exploiter des événements exportés dans une représentation
intermédiaire lorsque celle-ci est disponible.

Les champs manipulés dépendent du contenu effectivement fourni par l'événement :
horodatage, identifiant, fournisseur, canal, niveau et données associées.

L'objectif n'est pas de convertir artificiellement tous les événements Windows en une
ligne syslog, mais d'extraire un noyau commun sans supprimer les informations nécessaires
aux traitements spécifiques.

\subsection{HDFS et BGL}

HDFS et BGL sont des journaux système publics utilisés dans les expériences de séquences
et de templates.

Pour le protocole expérimental strict, Drain3 est utilisé dans un environnement séparé et
contrôlé. L'état du parseur est construit uniquement sur la partie d'entraînement avant
d'être gelé pour les ensembles de validation et de test.

Il faut distinguer ce protocole expérimental du pipeline multiformat général.

Cette distinction est importante parce que la validation finale montre que les adaptateurs
HDFS et BGL du pipeline multiformat ne produisent pas encore la couverture attendue. Leur
présence dans le code ne doit donc pas être interprétée comme une prise en charge complète
et universelle de ces formats.

\section{Normalisation}

\subsection{Implémentation du noyau commun}

Après parsing, le normaliseur tente de construire une représentation commune.

Les champs transversaux peuvent notamment concerner :

\begin{itemize}
    \item l'origine ;
    \item l'horodatage ;
    \item la famille ;
    \item le message ;
    \item la sévérité lorsqu'elle est connue ;
    \item l'hôte ou la source ;
    \item les informations nécessaires au routage.
\end{itemize}

Les informations propres à une famille restent accessibles lorsque leur conservation est
nécessaire au traitement.

Cette organisation évite de multiplier des interfaces entièrement différentes entre
chaque parser et les composants transversaux.

\subsection{Absence de garantie universelle}

L'implémentation ne permet pas d'affirmer que toute source traversant le système est
normalisée sans perte.

La campagne finale multiformat a précisément été conçue pour mesurer cette couverture.
Elle montre que certaines voies fonctionnent tandis que d'autres restent incomplètes.

Par conséquent, le code doit être décrit comme un ensemble d'adaptateurs disponibles et
non comme une fonction universelle capable de transformer n'importe quel log.

\subsection{Conservation du brut}

L'objectif architectural est de préserver les informations nécessaires à l'audit.
Cependant, la conservation exacte du message brut n'est pas universellement démontrée.

Il existe notamment des cas où la représentation d'origine n'est pas reproduite bit à bit
dans la sortie normalisée.

La formulation retenue dans le code et dans le mémoire distingue donc :

\begin{itemize}
    \item conservation d'informations exploitables ;
    \item présence d'un champ de message ;
    \item préservation exacte du brut.
\end{itemize}

Ces trois propriétés ne sont pas équivalentes.

\section{Implémentation du routeur}

\subsection{Fonction \texttt{route\_model}}

Le routeur réel est implémenté dans la logique associée à
\texttt{model\_router.py}. La fonction de routage examine les informations disponibles
sur la source et attribue des scores aux familles candidates.

Le processus peut tenir compte :

\begin{itemize}
    \item du nom et de l'extension ;
    \item de la structure du fichier ;
    \item des colonnes ;
    \item de motifs caractéristiques ;
    \item de champs reconnus dans les premières unités ;
    \item de l'existence d'un modèle compatible.
\end{itemize}

La décision finale contient suffisamment d'information pour être auditée.

\subsection{Sortie du routeur}

Une sortie de routage peut contenir notamment :

\begin{itemize}
    \item la famille proposée ;
    \item le modèle sélectionné ;
    \item les scores calculés ;
    \item les raisons ayant contribué au choix ;
    \item la marge entre les meilleurs candidats ;
    \item l'indication d'un fallback ;
    \item une erreur éventuelle ;
    \item la disponibilité du modèle.
\end{itemize}

Cette richesse est importante. Une simple étiquette finale ne permettrait pas de comprendre
pourquoi la décision a été prise.

\subsection{Interprétation de la confidence}

Le champ ou indicateur de \textit{confidence} utilisé par le routeur correspond à une
marge heuristique entre les deux meilleurs scores.

Il ne s'agit pas d'une probabilité calibrée.

On peut représenter cette marge de manière conceptuelle par :

\[
C_{\text{route}}
=
S_{(1)} - S_{(2)},
\]

où $S_{(1)}$ est le meilleur score heuristique et $S_{(2)}$ le deuxième.

La valeur indique donc à quel point le premier candidat se détache du second dans la
logique du routeur. Elle ne signifie pas qu'une famille a, par exemple, ``90\,\% de chance''
d'être correcte.

\subsection{Fallback}

Lorsque les indices sont insuffisants ou qu'aucun modèle compatible n'est disponible,
le routeur peut utiliser un fallback.

Ce comportement est volontaire. Il vaut mieux indiquer qu'aucune famille n'est
suffisamment compatible que forcer un modèle sur une entrée mal reconnue.

L'évaluation du routeur présentée au chapitre des résultats mesure directement ce
comportement, notamment sur le cas Apache qui conduit à une erreur de famille/fallback
dans le corpus utilisé.

\section{Registre des modèles}

\subsection{Rôle}

Le registre associe les familles de données aux traitements disponibles.

Il évite d'encoder directement dans tous les composants une liste de chemins d'artefacts
ou de modèles.

Un modèle exploitable doit être lié à suffisamment d'informations pour vérifier sa
compatibilité avec les données reçues.

Ces informations peuvent inclure :

\begin{itemize}
    \item famille ;
    \item type de modèle ;
    \item artefact ;
    \item variables attendues ;
    \item protocole ou version de référence ;
    \item informations d'intégrité disponibles.
\end{itemize}

\subsection{Séparation entre algorithme et artefact}

Dans le prototype, il faut distinguer :

\begin{enumerate}
    \item l'algorithme, par exemple Random Forest ;
    \item la configuration utilisée ;
    \item le modèle réellement entraîné ;
    \item le fichier qui contient cet artefact.
\end{enumerate}

Deux modèles Random Forest entraînés sur des données différentes ne doivent pas être
considérés comme le même détecteur simplement parce qu'ils utilisent le même algorithme.

Cette distinction devient particulièrement importante lors d'une mise à jour contrôlée.

\section{Implémentation des traitements de détection}

\subsection{Modèles supervisés}

Les traitements supervisés sont utilisés lorsque des labels fiables sont disponibles.

Dans les campagnes CICIDS2017, plusieurs candidats sont évalués :

\begin{itemize}
    \item Random Forest ;
    \item Extra Trees ;
    \item HistGradientBoosting ;
    \item Logistic Regression ;
    \item SGDClassifier avec fonction de perte logistique.
\end{itemize}

L'implémentation ne fixe pas l'un de ces modèles comme choix universel.

Les modèles sont comparés selon plusieurs métriques et plusieurs scénarios. Leur intérêt
est donc déterminé par les résultats expérimentaux plutôt que par une préférence
architecturale.

\subsection{Méthodes non supervisées et statistiques}

Les journaux ne possédant pas une vérité terrain complète nécessitent une autre approche.

Les méthodes utilisées dans les expériences HDFS/BGL comprennent notamment des traitements
statistiques et non supervisés.

Les sorties sont interprétées comme des scores ou anomalies candidates.

Aucune couche du code ne doit convertir automatiquement ces valeurs en attaque confirmée.

\subsection{Prétraitement}

Le prétraitement appartient au protocole de chaque modèle.

Pour les données tabulaires, cela peut comprendre :

\begin{itemize}
    \item sélection des colonnes ;
    \item transformation des valeurs non finies ;
    \item gestion des valeurs manquantes ;
    \item mise à l'échelle lorsque le modèle l'exige ;
    \item séparation des labels ;
    \item transformation dans le même espace de caractéristiques que celui de
    l'entraînement.
\end{itemize}

Une règle méthodologique importante est que toute transformation apprenant des paramètres
à partir des données doit être ajustée sur l'entraînement et non sur le test.

Cette contrainte appartient à la fois à l'implémentation et au protocole expérimental.

\section{Contrat métier des agents}

\subsection{Structure exacte}

Le contrat \texttt{AgentMessage} est défini dans
\texttt{src/logminer/agents/bus.py}.

Il contient exactement les sept champs suivants :

\begin{verbatim}
AgentMessage(
    run_id,
    source,
    target,
    message_type,
    payload,
    status,
    timestamp
)
\end{verbatim}

Le rôle de chaque champ est rappelé dans le tableau~\ref{tab:agentmessage-impl}.

\begin{table}[H]
\centering
\caption{Implémentation du contrat métier \texttt{AgentMessage}}
\label{tab:agentmessage-impl}
\small
\begin{tabularx}{\textwidth}{p{3cm}X}
\toprule
\textbf{Champ} & \textbf{Utilité} \\
\midrule

\texttt{run\_id} &
Corréler les échanges appartenant à une même exécution. \\

\texttt{source} &
Identifier le composant émetteur. \\

\texttt{target} &
Identifier le destinataire logique. \\

\texttt{message\_type} &
Indiquer la nature de la demande ou du résultat. \\

\texttt{payload} &
Transporter les informations spécifiques au traitement. \\

\texttt{status} &
Représenter l'état métier associé au message. \\

\texttt{timestamp} &
Associer un instant à l'échange. \\

\bottomrule
\end{tabularx}
\end{table}

Le contrat ne contient pas directement de champs
\texttt{event\_id}, \texttt{message\_id} ou \texttt{metadata}.

Des informations supplémentaires peuvent être incluses dans
\texttt{payload} lorsque le type du message le nécessite.

\subsection{Rôle de \texttt{run\_id}}

\texttt{run\_id} identifie le contexte d'une exécution et permet de rapprocher plusieurs
messages associés au même run.

Il ne doit pas être présenté comme l'identifiant unique de chaque événement de sécurité.

Cette distinction est nécessaire lorsque plusieurs agents échangent plusieurs messages
au cours d'une seule exécution.

\section{Bus local}

Le bus local permet aux composants de communiquer dans le même environnement d'exécution.

Son intérêt est de conserver le même contrat métier indépendamment de la topologie.

Un traitement peut donc être testé localement avant d'être confié à un worker utilisant
un mécanisme de transport plus explicite.

Cette séparation entre API métier et transport simplifie les tests unitaires et réduit le
couplage entre l'agent et Redis.

\section{Redis Streams}

\subsection{Utilisation}

Redis Streams est utilisé pour distribuer certaines tâches entre producteurs et workers.

Le mécanisme apporte notamment :

\begin{itemize}
    \item streams ;
    \item consumer groups ;
    \item consumers ;
    \item identifiants de transport ;
    \item messages pending ;
    \item acquittement ;
    \item mécanismes de reprise de certaines tâches.
\end{itemize}

Ces éléments appartiennent au niveau transport.

\subsection{Identifiant Redis et message métier}

Lorsqu'un message est publié dans Redis, le stream possède son propre identifiant.

Cet identifiant est utilisé pour des opérations telles que :

\begin{itemize}
    \item lecture ;
    \item suivi du pending ;
    \item acquittement ;
    \item reprise.
\end{itemize}

Il ne modifie pas la structure de \texttt{AgentMessage}.

Le code maintient donc deux niveaux :

\[
\text{message métier}
\neq
\text{enveloppe Redis}.
\]

Cette séparation est particulièrement importante dans l'interprétation de la campagne
multi-VM. Le fait qu'un message Redis soit lu et acquitté prouve le fonctionnement du
transport et du consumer group selon le protocole observé, mais ne constitue pas une
preuve générale que chaque conséquence applicative possible a réussi.

\subsection{Pending et récupération}

Lorsqu'un consumer récupère une tâche sans l'acquitter immédiatement, celle-ci peut rester
dans l'état pending.

Ce mécanisme permet de détecter certaines tâches abandonnées et, dans des scénarios
appropriés, de les confier à un autre worker.

L'implémentation fournit ainsi un support à la reprise.

Il ne faut cependant pas confondre cette propriété avec une garantie d'idempotence.
Le scénario où une tâche est effectivement traitée mais où le worker s'arrête juste avant
l'ACK peut éventuellement conduire à une nouvelle livraison. Les conséquences métier de
ce cas particulier ne sont pas évaluées dans les campagnes finales.

\section{MQTT}

Le prototype comprend également une voie de communication MQTT.

Le principe retenu reste identique : le contrat métier est sérialisé avant d'être envoyé
par le protocole de transport.

MQTT ne crée pas de nouveau champ métier dans \texttt{AgentMessage}. Les propriétés
techniques propres au protocole doivent rester séparées du contenu métier.

Dans le mémoire, MQTT est présenté comme un mécanisme implémenté de communication et non
comme la preuve d'une infrastructure distribuée de production.

\section{Supervision des agents}

\subsection{Capacités}

Un agent peut déclarer plusieurs capacités correspondant aux tâches qu'il sait exécuter.

Cette information permet d'éviter qu'un worker traite une tâche pour laquelle il ne
dispose ni du handler ni des dépendances nécessaires.

La sélection peut donc prendre en compte :

\begin{itemize}
    \item type de tâche ;
    \item capacités disponibles ;
    \item état de l'agent ;
    \item historique récent ;
    \item disponibilité.
\end{itemize}

\subsection{Heartbeat}

Les agents peuvent émettre un heartbeat permettant au superviseur de disposer d'une
indication sur leur activité.

Le heartbeat ne constitue pas à lui seul un mécanisme complet de haute disponibilité.

Il permet principalement de détecter qu'un agent ne se manifeste plus selon le comportement
attendu et de soutenir certaines décisions d'orchestration.

\subsection{Mémoire locale du superviseur}

Le superviseur conserve des informations sur les exécutions et les choix récents.

La mémoire peut notamment contribuer à éviter de sélectionner systématiquement les mêmes
sources lorsqu'une politique de rotation ou de priorité est appliquée.

L'intérêt causal de cette mémoire sur la qualité de priorisation n'est cependant pas
isolé expérimentalement dans le travail final.

Il est donc correct d'affirmer que le mécanisme existe et influence certaines décisions
du code, mais pas qu'il améliore quantitativement la priorisation dans tous les cas.

\section{Mémoire persistante et audit}

\subsection{Deux fonctions complémentaires}

Mémoire et audit sont proches mais ne remplissent pas exactement le même rôle.

La mémoire sert à conserver des informations réutilisables lors de traitements ultérieurs.

L'audit cherche à reconstruire ce qui s'est produit pendant une exécution.

Un même événement peut donc contribuer aux deux mécanismes.

\subsection{Informations conservées}

Selon les modules et les campagnes, les informations persistantes peuvent comprendre :

\begin{itemize}
    \item identifiant d'exécution ;
    \item source ;
    \item agent ;
    \item type d'action ;
    \item statut ;
    \item erreur ;
    \item résultat ;
    \item décision analyste ;
    \item famille ;
    \item modèle ;
    \item instant ;
    \item informations d'intégrité.
\end{itemize}

La persistance permet de relier l'application aux campagnes expérimentales.

\subsection{Décisions analyste}

Lorsqu'une fonction de validation est utilisée, une anomalie candidate peut être
acceptée, rejetée ou reclassée.

Cette décision est enregistrée afin de rester disponible pour l'audit et les traitements
ultérieurs.

Elle ne provoque pas directement un réentraînement automatique du modèle actif.

Cette séparation constitue une mesure de sécurité méthodologique : une erreur humaine
isolée ne doit pas immédiatement modifier le comportement du détecteur.

\section{Implémentation de la corrélation}

\subsection{Entrée}

Le corrélateur reçoit des anomalies candidates auxquelles sont associés des attributs
contextuels.

Toutes les sorties de modèles ne disposent pas des mêmes champs. L'implémentation doit donc
utiliser uniquement les informations réellement disponibles.

\subsection{Logique de regroupement}

La corrélation repose sur des règles combinant notamment :

\begin{itemize}
    \item proximité temporelle ;
    \item source ;
    \item famille ;
    \item identifiants ou attributs communs ;
    \item autres informations disponibles dans les événements.
\end{itemize}

Le moteur produit des groupes correspondant à des incidents candidats.

\subsection{Traçabilité du regroupement}

Un regroupement doit rester explicable.

L'audit doit permettre de retrouver les événements ayant contribué à un incident candidat
et les critères utilisés pour les rapprocher.

Cette propriété est importante dans l'expérience synthétique finale, où la qualité du
regroupement est évaluée indépendamment de la détection en amont.

\section{Mise à jour contrôlée des modèles}

\subsection{Principe d'implémentation}

La mise à jour ne remplace pas directement le modèle actif après un nouvel entraînement.

Le code distingue :

\begin{itemize}
    \item modèle courant ;
    \item candidat ;
    \item score courant ;
    \item score candidat ;
    \item différence ;
    \item seuil minimal ;
    \item décision.
\end{itemize}

La différence est définie par :

\[
\Delta M
=
M_{\text{candidat}}
-
M_{\text{courant}}.
\]

La promotion est autorisée lorsque :

\[
\Delta M \geq \delta_{\min}
\]

pour une métrique où une valeur plus élevée représente une amélioration.

\subsection{Sauvegarde et intégrité}

Avant une promotion, l'ancien modèle doit rester identifiable.

Des hashes permettent de vérifier l'intégrité des artefacts utilisés dans la procédure.

Le fonctionnement final démontre trois branches :

\begin{itemize}
    \item promotion ;
    \item rejet pour dégradation ;
    \item rejet pour amélioration positive mais inférieure au seuil.
\end{itemize}

Les valeurs numériques correspondantes seront présentées dans le chapitre expérimental.

\subsection{Portée}

Cette implémentation démontre un mécanisme de mise à jour contrôlée de bout en bout.

Elle ne démontre pas :

\begin{itemize}
    \item un service autonome de réentraînement permanent ;
    \item un déploiement automatique en production ;
    \item une détection automatique du concept drift ;
    \item une amélioration continue garantie.
\end{itemize}

Ces fonctions restent des extensions possibles.

\section{API locale}

\subsection{Rôle}

L'API locale sépare le moteur d'analyse de l'interface utilisateur.

Cette couche évite que le dashboard dépende directement de l'organisation interne des
fichiers ou des modules Python.

Elle peut exposer des fonctions permettant notamment de :

\begin{itemize}
    \item lancer ou consulter une analyse ;
    \item accéder aux résultats ;
    \item consulter certaines ressources ;
    \item récupérer des informations d'audit ;
    \item enregistrer une décision analyste ;
    \item vérifier l'état de certains services.
\end{itemize}

\subsection{Intérêt architectural}

Cette séparation apporte trois bénéfices.

Premièrement, l'interface peut évoluer sans modifier directement les modèles.

Deuxièmement, les fonctions exposées peuvent être testées indépendamment de l'affichage.

Troisièmement, l'API constitue une frontière explicite entre les données présentées à
l'utilisateur et le stockage interne.

Le mémoire ne présente toutefois pas cette API comme un service Internet de production.
Elle appartient au prototype local.

\section{Dashboard Ariel Logminer}

\subsection{Interface finale retenue}

Le prototype possède une interface web destinée à la consultation des résultats.

Dans la version finale du mémoire, seule \textbf{l'interface à dominante bleue} est
retenue comme interface de référence.

Les anciennes captures à dominante rouge-or ne font pas partie de la présentation finale.

Le dashboard bleu doit être utilisé pour les futures figures du mémoire.

\subsection{Fonctions principales}

L'interface organise plusieurs catégories d'information :

\begin{itemize}
    \item indicateurs généraux ;
    \item anomalies candidates ;
    \item incidents candidats ;
    \item détails associés aux résultats ;
    \item filtres ;
    \item ressources ou état de certains composants ;
    \item informations d'audit ;
    \item décisions analyste lorsque la fonction est disponible.
\end{itemize}

L'objectif est de permettre au lecteur ou à l'analyste de passer d'une vue synthétique à
un détail exploitable.

L'interface ne remplace pas les données expérimentales. Les CSV, JSON et rapports restent
les sources utilisées pour les calculs scientifiques ; le dashboard constitue une couche
de consultation.

\subsection{Validation fonctionnelle}

Les fonctions visibles du dashboard permettent d'établir que l'interface est intégrée au
prototype et qu'elle peut présenter certaines sorties de la chaîne.

Cette preuve est de nature fonctionnelle.

Elle ne permet pas de conclure que :

\begin{itemize}
    \item l'interface réduit le temps d'analyse ;
    \item elle diminue la charge cognitive ;
    \item elle est plus ergonomique qu'une autre interface ;
    \item les analystes la préfèrent ;
    \item elle améliore la qualité des décisions.
\end{itemize}

Ces affirmations nécessiteraient une étude utilisateur spécifique.

% ============================================================
% DASHBOARD BLEU
% ============================================================
% NE PAS réactiver les anciennes captures rouge-or.
%
% Lorsque les vraies captures de l'interface bleue seront disponibles,
% les placer par exemple dans :
%
% figures/dashboard/dashboard_bleu_vue_ensemble.png
% figures/dashboard/dashboard_bleu_resultats.png
% figures/dashboard/dashboard_bleu_detail_incident.png
% figures/dashboard/dashboard_bleu_audit.png
%
% Puis intégrer uniquement les captures bleues.
% ============================================================

\section{Formats de stockage et artefacts}

\subsection{CSV}

CSV est utilisé lorsque les données possèdent une structure tabulaire.

Ce format présente plusieurs avantages pour les expériences :

\begin{itemize}
    \item lecture simple ;
    \item compatibilité avec pandas ;
    \item inspection manuelle ;
    \item importation dans de nombreux outils ;
    \item génération directe de tableaux.
\end{itemize}

Il est particulièrement adapté aux résultats de campagnes contenant une ligne par run,
seed, scénario ou modèle.

\subsection{JSON et JSONL}

JSON est utilisé pour des rapports plus structurés, par exemple lorsque plusieurs
informations hiérarchiques doivent être conservées dans un même artefact.

JSONL permet de stocker une succession d'objets indépendants, un par ligne.

Ces formats sont adaptés aux journaux d'audit et aux sorties de traitements où toutes les
observations ne possèdent pas nécessairement exactement la même structure tabulaire.

\subsection{Artefacts de modèles}

Les modèles entraînés sont conservés séparément des résultats.

Cette séparation permet d'utiliser les hashes pour identifier l'artefact actif et de
distinguer :

\begin{itemize}
    \item données ;
    \item modèle ;
    \item métriques ;
    \item rapport.
\end{itemize}

\section{Gestion des erreurs}

\subsection{Principe}

Une erreur est considérée comme un résultat technique à enregistrer.

Une exception masquée ou un fichier silencieusement ignoré rend l'analyse scientifique
difficile, car il devient impossible de connaître le nombre réel d'entrées traitées.

Le système cherche donc à conserver une information explicite sur :

\begin{itemize}
    \item l'entrée concernée ;
    \item le composant ;
    \item le type d'échec ;
    \item le statut ;
    \item le message disponible.
\end{itemize}

\subsection{Erreurs de parsing}

Une erreur de parsing peut conduire :

\begin{itemize}
    \item à un fallback ;
    \item à une sortie partielle ;
    \item à un rejet explicite.
\end{itemize}

Le choix dépend de la quantité d'information encore exploitable.

\subsection{Erreurs de modèle}

Un modèle peut être indisponible ou incompatible avec les colonnes reçues.

Dans ce cas, le système ne doit pas générer artificiellement une prédiction.

Le routeur et le registre sont précisément destinés à réduire ce risque.

\subsection{Erreurs de transport}

Une erreur Redis possède une nature différente d'une erreur de détection.

Le message peut ne pas avoir été livré, être pending ou avoir été acquitté.

L'audit doit donc conserver une distinction entre l'état du transport et celui du traitement
métier.

\section{Portabilité}

\subsection{Dépendances principales}

Le prototype cherche à maintenir un noyau relativement léger.

Python, pandas, NumPy et scikit-learn constituent la base de nombreuses opérations.

Certaines fonctions dépendent de bibliothèques ou services complémentaires tels que
Drain3 ou Redis.

L'absence d'une dépendance optionnelle ne doit pas rendre toutes les fonctions du
prototype inutilisables.

\subsection{Systèmes d'exploitation}

Les campagnes finales sont principalement exécutées dans un environnement Windows.

Certaines validations utilisent également Debian et Ubuntu dans des machines virtuelles.

Le code Python générique peut fonctionner sur plusieurs systèmes, mais les scripts
d'administration ou de collecte liés à Windows ne sont pas automatiquement portables sous
Linux.

Cette distinction doit être conservée dans la documentation de reproduction.

\section{Sécurité de l'implémentation}

Le prototype analyse potentiellement des informations sensibles.

Plusieurs précautions de conception sont donc pertinentes :

\begin{itemize}
    \item limiter la copie inutile de données ;
    \item utiliser des chemins locaux lorsque cela suffit ;
    \item ne pas exposer automatiquement toutes les informations brutes via l'API ;
    \item anonymiser ou neutraliser les informations sensibles dans les artefacts destinés
    à la publication ;
    \item séparer les données expérimentales des captures et figures publiques.
\end{itemize}

Ces précautions ne constituent pas une certification de cybersécurité du prototype.

Un déploiement de production nécessiterait notamment une analyse spécifique concernant :

\begin{itemize}
    \item authentification ;
    \item contrôle d'accès ;
    \item chiffrement ;
    \item gestion des secrets ;
    \item durcissement Redis ;
    \item exposition réseau de l'API ;
    \item politique de rétention.
\end{itemize}

Ces aspects sont volontairement distingués des fonctions évaluées dans le mémoire.

\section{Lien entre implémentation et évaluation}

Le tableau~\ref{tab:implementation-evaluation} relie les principaux composants du code aux
preuves utilisées dans la suite du mémoire.

\begin{table}[H]
\centering
\caption{Lien entre les composants implémentés et leur mode d'évaluation}
\label{tab:implementation-evaluation}
\small
\begin{tabularx}{\textwidth}{p{3.5cm}p{4cm}X}
\toprule
\textbf{Composant} &
\textbf{Type de preuve} &
\textbf{Interprétation correcte} \\
\midrule

Pipeline principal &
Campagnes fonctionnelles sur plusieurs voies. &
Le pipeline est implémenté et fonctionne pour les voies effectivement observées. \\

Normalisation &
Validation multiformat. &
Couverture partielle ; aucune robustesse universelle n'est revendiquée. \\

Routeur &
Corpus dérivé de 81 fichiers. &
Évaluation fonctionnelle dans ce corpus ; pas de généralisation universelle. \\

Détecteurs &
Métriques CICIDS2017, HDFS et BGL. &
Performance liée au dataset et au protocole. \\

Agents &
Benchmark monolithe/agents et campagnes Redis. &
Modularité et distribution observées ; pas de gain de débit systématique. \\

Redis Streams &
Lecture, ACK, pending et reprise de laboratoire. &
Transport distribué démontré dans le périmètre testé. \\

Corrélateur &
Scénarios synthétiques contrôlés. &
Qualité mesurable sur ces scénarios ; pas de validation SOC réelle. \\

Mise à jour &
Trois branches contrôlées. &
Promotion/rejet fonctionnels ; pas d'apprentissage continu autonome. \\

Dashboard &
Inspection fonctionnelle des vues. &
Interface intégrée ; utilisabilité non évaluée. \\

\bottomrule
\end{tabularx}
\end{table}

Cette correspondance entre composant et preuve évite qu'une simple présence dans le code
soit interprétée comme une validation expérimentale complète.

\section{Statut des principales fonctions implémentées}

Le tableau~\ref{tab:statut-implementation} synthétise le niveau réellement atteint par les
fonctions principales.

\begin{table}[H]
\centering
\caption{Statut des principales fonctions du prototype}
\label{tab:statut-implementation}
\small
\begin{tabularx}{\textwidth}{p{4.2cm}X}
\toprule
\textbf{Élément} & \textbf{Statut retenu} \\
\midrule

Pipeline principal &
Implémenté et testé fonctionnellement selon plusieurs voies. \\

\texttt{AgentMessage} &
Implémenté avec sept champs métier exactement. \\

Bus local &
Implémenté. \\

Redis Streams &
Implémenté et testé dans des campagnes locales et multi-VM de laboratoire. \\

MQTT &
Implémenté comme voie de communication ; ne constitue pas une preuve de distribution
industrielle. \\

Routeur &
Implémenté et évalué sur un corpus dérivé. \\

Corrélateur &
Implémenté et évalué sur des scénarios synthétiques contrôlés. \\

Mise à jour contrôlée &
Implémentée et testée fonctionnellement de bout en bout. \\

Apprentissage continu autonome &
Non démontré ; perspective. \\

Adaptateurs multiformat &
Fonctionnels pour plusieurs sources mais couverture incomplète, notamment dans les voies
HDFS/BGL de la campagne multiformat. \\

Reprise avant traitement/ACK &
Testée dans un scénario de laboratoire. \\

Idempotence après traitement avant ACK &
Non évaluée. \\

Haute disponibilité de production &
Non évaluée. \\

\bottomrule
\end{tabularx}
\end{table}

\section{Choix de légèreté}

Le prototype privilégie des méthodes et outils compatibles avec un environnement local.

Ce choix présente plusieurs intérêts.

Premièrement, il réduit les ressources nécessaires pour reproduire une partie importante
des expériences.

Deuxièmement, il facilite l'inspection des modèles et des données intermédiaires.

Troisièmement, il permet de concentrer le travail sur l'intégration et les protocoles
plutôt que sur l'entraînement de modèles très lourds.

Cette décision ne signifie pas que les architectures profondes seraient inadaptées à
l'analyse de logs.

Comme l'a montré l'état de l'art, les modèles séquentiels et Transformers peuvent être très
pertinents lorsque l'objectif est de représenter des dépendances complexes entre
événements.

Dans Ariel \logminer{}, ces approches constituent principalement des perspectives ou des
points de comparaison conceptuels. Le prototype final repose surtout sur des modèles et
méthodes dont l'entraînement et l'évaluation peuvent être contrôlés dans les ressources
disponibles.

\section{Implémentation et reproductibilité}

L'un des objectifs de l'implémentation est de maintenir une relation directe entre le code
et les résultats scientifiques.

Cela implique notamment :

\begin{itemize}
    \item conserver les paramètres des campagnes ;
    \item utiliser des seeds explicites lorsqu'une méthode est stochastique ;
    \item enregistrer les sorties dans des fichiers identifiables ;
    \item calculer les hashes des données ou modèles lorsque nécessaire ;
    \item séparer les fichiers utilisés pour l'entraînement, la validation et le test ;
    \item conserver les échecs et résultats négatifs ;
    \item produire des figures à partir des résultats consolidés plutôt que de valeurs
    recopiées manuellement.
\end{itemize}

Le chapitre consacré à la reproductibilité détaillera l'environnement, les scripts et les
artefacts finaux.

À ce stade, le point essentiel est que la reproductibilité n'a pas été ajoutée uniquement
après les expériences. Elle influence directement la manière dont les modules, sorties et
campagnes sont organisés.

\section{Limites de l'implémentation actuelle}

L'implémentation finale atteint le niveau d'un prototype fonctionnel, mais plusieurs limites
restent importantes.

La première concerne la couverture multiformat. Plusieurs adaptateurs produisent des sorties
exploitables, tandis que certaines voies restent incomplètes. La présence d'un parser ou
d'une fonction dans le dépôt ne doit donc pas être assimilée à une compatibilité complète.

La deuxième concerne la distribution. Redis Streams et plusieurs workers permettent une
validation multi-VM, mais l'infrastructure ne comprend pas de mécanisme complet de
réplication et de basculement démontré.

La troisième concerne la reprise. Une tâche pending peut être reprise dans le scénario
testé, mais l'idempotence métier après traitement et avant ACK n'est pas établie.

La quatrième concerne la mémoire. Les informations historiques existent et peuvent
contribuer à certaines décisions du superviseur, mais leur bénéfice causal n'a pas été
isolé expérimentalement.

La cinquième concerne le dashboard. Son fonctionnement peut être observé, mais aucune
étude utilisateur ne mesure son ergonomie ou son effet sur la qualité de l'analyse.

La sixième concerne l'adaptation. La procédure de promotion d'un modèle candidat est
fonctionnelle, mais elle ne constitue pas un système autonome d'apprentissage continu en
production.

Ces limites ne remettent pas en cause l'existence des composants. Elles déterminent le
niveau exact de revendication que l'implémentation permet de soutenir.

\section{Synthèse du chapitre}

Ce chapitre a présenté la traduction logicielle de l'architecture d'Ariel \logminer{}.

Le pipeline principal organise le parcours d'une source depuis sa lecture jusqu'au
traitement analytique. Les parseurs et adaptateurs transforment des représentations
hétérogènes, tandis que le normaliseur fournit un noyau commun lorsque la voie le permet.
Le routeur examine les caractéristiques de la source et sélectionne une famille de
traitement ; sa confidence reste une marge heuristique et non une probabilité calibrée.

Le portefeuille de modèles permet d'associer plusieurs traitements aux familles disponibles
sans supposer l'existence d'un détecteur universel. Les sorties produites sont conservées
comme classes, scores ou anomalies candidates selon la nature de la méthode.

La couche multi-agent utilise un contrat métier minimal composé exactement de
\texttt{run\_id}, \texttt{source}, \texttt{target},
\texttt{message\_type}, \texttt{payload}, \texttt{status} et
\texttt{timestamp}. Les identifiants de Redis appartiennent au transport et restent
séparés de ce contrat.

Redis Streams apporte des mécanismes de distribution, d'acquittement et de reprise de
certaines tâches. Le superviseur utilise les capacités et états des agents pour organiser
leur exécution. La mémoire et l'audit conservent les informations nécessaires à la
traçabilité, tandis que le corrélateur rapproche plusieurs anomalies candidates.

La mise à jour des modèles repose sur une logique conservatrice : un candidat est comparé
au modèle courant avant toute promotion. Le prototype démontre donc une mise à jour
contrôlée, mais pas un apprentissage continu autonome en production.

Enfin, le dashboard constitue la couche d'interaction avec l'analyste. Dans la version
finale du mémoire, seule l'interface bleue doit être utilisée. Son intégration démontre une
fonction de consultation, mais aucune conclusion d'utilisabilité n'est formulée sans étude
utilisateur.

Ces choix rendent le prototype suffisamment concret pour être soumis à plusieurs formes
d'évaluation. Le chapitre suivant présente la méthodologie expérimentale finale et les
résultats obtenus sur CICIDS2017, HDFS, BGL, le pipeline multiformat, le routeur, les
agents, Redis, la corrélation et la mise à jour contrôlée.